% ---------------------------------------------------------------------------
% Table II -- the backbones. Generated by experiments/make_backbone_table.py.
% Regenerate rather than edit: the layer counts and the actions-per-call
% column are derived from the run records, and the other columns assert
% against the run scripts they were copied from.
%
% Spans both columns, so table* and not table. Needs \usepackage{booktabs}.
%
% This table exists to keep four facts out of the prose. Setup must not also
% say in words that SpatialVLA has 26 layers, that UniVLA emits five actions
% a call, or which card ran what. It may say what those facts MEAN, which is
% a different sentence: four removed layers is 12.5% of one stack and 15.4%
% of another, and action repeat 4 is a 4-step horizon for two backbones and
% a 20-step one for the third.
%
% "Actions per call" and not "native chunk size". Native chunk is a property
% of a checkpoint; this column is a property of our runs, and for SpatialVLA
% the two need not agree, since the SimplerEnv path keeps raw_actions[0] and
% discards the rest. The statistics multiply THIS number by the action
% repeat, so this is the one the reader needs.
%
% The params column is not a specification for its own sake. It is why the
% GPU column has two values in it: 8.5B at fp16 is 17 GB and a T4 holds
% about 15.
%
% The UniVLA folder name is 37 characters and sits in the note rather than
% in the column, which is a width decision and not an editorial one. At 10pt
% typewriter a character is about 5.25pt, so putting the repository and the
% folder in one cell would set that cell at roughly 285pt of the 505.89pt
% \textwidth ieeeconf gives, and the remaining five columns need about 240.
% With the folder in the note the widest cell is the 35-character SpatialVLA
% identifier at about 185pt, and the row totals roughly 460pt including
% column separation. If a column is added, redo this arithmetic.
% ---------------------------------------------------------------------------
\begin{table*}[t]
\centering
\caption{The three backbones, their public checkpoints, and how each was run.}
\label{tab:backbones}
\setlength{\tabcolsep}{6pt}
\begin{tabular}{l l l r r l}
\toprule
Backbone & Checkpoint & Params & Decoder layers & Actions per call & GPU \\
\midrule
OpenVLA & \texttt{openvla/openvla-7b} & 7B & $32$ & $1$ & T4 (15\,GB) \\
SpatialVLA & \texttt{IPEC-COMMUNITY/spatialvla-4b-224-pt} & 4B & $26$ & $1$ & T4 (15\,GB) \\
UniVLA & \texttt{Yuqi1997/UniVLA} & 8.5B & $32$ & $5$ & L4 (22.5\,GB) \\
\bottomrule
\end{tabular}

\parbox{\textwidth}{\footnotesize Actions per call is how many actions one
forward pass contributes to the environment, so an action repeat of $r$ puts
$r$ times that many environment steps between decisions. Layer counts are
recovered from the pruning records rather than read from the configuration
files. The \texttt{Yuqi1997/UniVLA} weights are the \texttt{UNIVLA\_SIMPLER\_BRIDGE\_VIDEO\_BS128\_20K} folder of that repository.}
\end{table*}
